% !TeX program = lualatex
% !TeX encoding = utf8
% !TeX spellcheck = uk_UA
% !TeX root = ../main.tex

%% --------------------------------------------------------
%% Chapter file: Waves.tex
%% --------------------------------------------------------


\chapter{Колективні коливання та хвилі в~плазмі}\label{sec:plasma-waves}\hypertarget{plasma-waves}{}
\graphicspath{{\currfilebase/Pictures}}
%% --------------------------------------------------------

Дебаївська довжина $\lambda_D$ задає просторовий масштаб колективної відповіді плазми. Тут
ми введемо її часовий аналог~--- плазмову частоту $\omega_p$~--- і~розглянемо колективні
збурення, що змінюються і~в~просторі, і~в~часі. Тут з'являються два
різні класи хвиль: поперечні електромагнітні хвилі та поздовжні електростатичні плазмові
хвилі. Їх об'єднує спільне джерело~--- колективний рух заряджених частинок, але фізика
поширення в~них різна.



\section{Плазмові коливання та плазмова частота}\label{sec:plasma-oscillations}
%% --------------------------------------------------------



Дебаївське екранування описує, як плазма просторово перебудовується навколо внесеного
заряду. Розглянемо тепер, як плазма реагує на збурення, що змінюється з~часом. Нехай
електронний компонент однорідної плазми змістився відносно нерухомого іонного фону на
невелику відстань $x$ (рис.~\ref{pic:plasma-oscillations}). Іони можна вважати нерухомими,
оскільки $m_i\gg m_e$.

На краях шару виникають заряди з~поверхневою густиною $\sigma=\mp n_eex$, і~між ними~---
поле плоского конденсатора, напрямлене так, щоб повернути електрони назад:
% ~~~~~~~~~~~~~~~~~~~~~~~~~~~~~~~~~~~~~~~~~~~~~~~~~~~~~~~~
\begin{equation*}
	E=4\pi n_eex.
\end{equation*}
% ~~~~~~~~~~~~~~~~~~~~~~~~~~~~~~~~~~~~~~~~~~~~~~~~~~~~~~~~
Рівняння руху електрона (його заряд $-e$)
% ~~~~~~~~~~~~~~~~~~~~~~~~~~~~~~~~~~~~~~~~~~~~~~~~~~~~~~~~
\begin{equation}\label{eq:plasma-oscillation-eq}
	m_e\ddot x=-eE=-4\pi n_ee^2x
	\quad\Rightarrow\quad
	\ddot x+\omega_{pe}^2x=0,
\end{equation}
% ~~~~~~~~~~~~~~~~~~~~~~~~~~~~~~~~~~~~~~~~~~~~~~~~~~~~~~~~
де
% ~~~~~~~~~~~~~~~~~~~~~~~~~~~~~~~~~~~~~~~~~~~~~~~~~~~~~~~~
\begin{equation}\label{eq:plasma-frequency}
	\tcbhighmath{\omega_{pe}=\sqrt{\frac{4\pi n_ee^2}{m_e}}\approx5{,}64\times10^4\sqrt{n_e[\mathrm{см}^{-3}]}~\mathrm{с}^{-1}}
\end{equation}
% ~~~~~~~~~~~~~~~~~~~~~~~~~~~~~~~~~~~~~~~~~~~~~~~~~~~~~~~~
називається \emphx{електронною плазмовою} (або \emphx{ленгмюрівською}) \emphx{частотою}.
Електрони гармонічно коливаються відносно іонів з~цією частотою, і~ці коливання називають
\emphx{плазмовими}.

\begin{figure}[h!]\centering
    %---------------------------------------------------------
	\begin{subfigure}[t]{0.5\linewidth}\centering
		\localinput{plasma_ions.tikz}
	\end{subfigure}%
	%---------------------------------------------------------
	\begin{subfigure}[t]{0.5\linewidth}\centering
		\localinput{condensator.tikz}
	\end{subfigure}
	%---------------------------------------------------------
	\caption{Модель плазмових коливань: електронний компонент зміщений відносно нерухомого
	позитивного фону, виникає повертаюче електричне поле}
	\label{pic:plasma-oscillations}
\end{figure}

Той самий результат дає й~опис без уявного конденсатора. З~рівняння неперервності
$\partial_t\rho+\nabla\cdot\vect{j}=0$ та~рівняння руху $m_e\partial_t\vect{v}=-e\vect{E}$
при $\vect{j}=-n_ee\vect{v}$ маємо
$\partial_t^2\rho=-\nabla\cdot\partial_t\vect{j}=-(n_ee^2/m_e)\nabla\cdot\vect{E}$, а~за
теоремою Гаусса $\nabla\cdot\vect{E}=4\pi\rho$. Тому
% ~~~~~~~~~~~~~~~~~~~~~~~~~~~~~~~~~~~~~~~~~~~~~~~~~~~~~~~~
\begin{equation*}
	\frac{\partial^2\rho}{\partial t^2}+\omega_{pe}^2\rho=0.
\end{equation*}
% ~~~~~~~~~~~~~~~~~~~~~~~~~~~~~~~~~~~~~~~~~~~~~~~~~~~~~~~~
Тобто плазмова частота~--- це частота, з~якою плазма \enquote{відновлює} порушену
квазінейтральність, а~$\omega_{pe}^{-1}$~--- характерний час цього відновлення.

Якщо врахувати рух іонів (заряд $Ze$, маса $m_i$), для холодної двокомпонентної плазми
повна частота довгохвильових поздовжніх коливань
% ~~~~~~~~~~~~~~~~~~~~~~~~~~~~~~~~~~~~~~~~~~~~~~~~~~~~~~~~
\begin{equation*}
	\omega_p^2=\omega_{pe}^2+\omega_{pi}^2=4\pi e^2\left(\frac{n_e}{m_e}+\frac{Z^2n_i}{m_i}\right)\simeq\omega_{pe}^2.
\end{equation*}
% ~~~~~~~~~~~~~~~~~~~~~~~~~~~~~~~~~~~~~~~~~~~~~~~~~~~~~~~~~~~~~~
Іонний внесок мізерний (у~задачі~\ref{prb:PlasmaIons} його необхідно оцінити). Крім того,
коливання в~реальній плазмі загасають (зіткнення, теплове розмазування), а~при врахуванні
теплового руху їхня частота залежить від хвильового числа (§\ref{sec:langmuir-waves}).
Тому $\omega_{pe}$ слід розуміти насамперед як базову \emph{масштабну} частоту плазми.



\subsection*{Зв'язок двох масштабів}



Ми вже знаємо просторовий масштаб відповіді плазми. Тепер поставимо природне наступне
питання: \emphx{як швидко} плазма відновлює порушену квазінейтральність? Дебаївську
довжину та плазмову частоту поєднує проста формула
% ~~~~~~~~~~~~~~~~~~~~~~~~~~~~~~~~~~~~~~~~~~~~~~~~~~~~~~~~
\begin{equation}\label{eq:debye-vT-wp}
	\tcbhighmath{\lambda_{De}=\frac{v_{Te}}{\omega_{pe}},\qquad v_{Te}=\sqrt{\frac{k_{\mathrm B}T_e}{m_e}}\approx4{,}2\times10^{7}\sqrt{T_e[\text{еВ}]}~\mathrm{см/с}}.
\end{equation}
% ~~~~~~~~~~~~~~~~~~~~~~~~~~~~~~~~~~~~~~~~~~~~~~~~~~~~~~~~~~~~~~
Дебаївська довжина~--- це відстань, яку електрон із~тепловою швидкістю проходить за час
порядку $\omega_{pe}^{-1}$. Одразу видно, що плазма на масштабах менше $\lambda_D$
\enquote{не встигає} колективно відповідати~--- тепловий рух розмиває ефект. Фізичний
зміст двох масштабів:
% ~~~~~~~~~~~~~~~~~~~~~~~~~~~~~~~~~~~~~~~~~~~~~~~~~~~~~~~~
\begin{enumerate}
	\item час колективної реакції $\sim\omega_{pe}^{-1}$;
	\item довжина колективної реакції $\sim\lambda_{De}$.
\end{enumerate}
% ~~~~~~~~~~~~~~~~~~~~~~~~~~~~~~~~~~~~~~~~~~~~~~~~~~~~~~~~~~~~~~~~
Для швидких оцінок зручна формула
$\lambda_{De}\approx7{,}43\times10^{2}\sqrt{T_e[\text{еВ}]/n_e[\mathrm{см}^{-3}]}$~см.

Отже, в~плазмі виникає принципова ієрархія просторових масштабів:
% ~~~~~~~~~~~~~~~~~~~~~~~~~~~~~~~~~~~~~~~~~~~~~~~~~~~~~~~~
\begin{equation}\label{eq:scale-hierarchy}
	a\ll\lambda_D\ll L.
\end{equation}
% ~~~~~~~~~~~~~~~~~~~~~~~~~~~~~~~~~~~~~~~~~~~~~~~~~~~~~~~~
Перша нерівність~--- умова $N_D\gg1$ (екранування колективне й~статистичне), друга~---
умова квазінейтральності в~об'ємі. У~діапазоні $a\ll r\ll\lambda_D$ ще важливі
мікроскопічні властивості окремих частинок, але екранування на масштабі $\lambda_D$ уже є
колективним. На масштабах $L\gg\lambda_D$ можна вводити гладкі макроскопічні поля та
густини.



%% --------------------------------------------------------
\section{Діелектрична проникність холодної плазми}\label{sec:plasma-permittivity}
%% --------------------------------------------------------



Нехай на плазму падає електромагнітна хвиля з~частотою $\omega$ (рис.~\ref{tikz:emwave}).
Знайдемо відповідь плазми на змінне поле. Зробимо чотири припущення, які в~типових умовах
виконуються.

%---------------------------------------------------------
\begin{figure}[h!]\centering
\localinput{emwave.tikz}
\caption{Електромагнітна хвиля, що падає з~вакууму на шар плазми}
\label{tikz:emwave}
\end{figure}
%---------------------------------------------------------
\begin{enumerate}
	\item Рухаються переважно електрони. Маса іона не менша за $1836\,m_e$, тому струм
	  іонів малий порівняно зі струмом електронів (внесок іонів пропорційний $m_e/m_i$).
	\item Швидкість електронів мала, $v\ll c$, тому силою Лоренца від магнітного поля
	  хвилі ($\sim(v/c)eE$) можна знехтувати.
	\item На масштабі руху електронів поле можна вважати однорідним. Для теплового руху це
	  означає $kv_{Te}\ll\omega$; амплітуда вимушеного руху також має бути малою, щоб
	  зберігалася лінеаризація.
	\item Зіткненнями можна знехтувати на часі одного коливання: $\nu\ll\omega$.
\end{enumerate}

Рівняння руху електрона $m_e\ddot{\vect r}=-e\vect{E}$ має розв'язок
% ~~~~~~~~~~~~~~~~~~~~~~~~~~~~~~~~~~~~~~~~~~~~~~~~~~~~~~~~
\begin{equation*}
	\vect r=\frac{e}{m_e\omega^2}\vect E,
\end{equation*}
% ~~~~~~~~~~~~~~~~~~~~~~~~~~~~~~~~~~~~~~~~~~~~~~~~~~~~~~~~
тобто зміщення збігається за напрямом з~полем $\vect E$. Це означає, що воно перебуває
в~протифазі до сили $-e\vect E$, яка діє на електрон (так завжди буває при вимушених
коливаннях вільної частинки на частоті, більшій за власну). Електрони плазми подібні до
зв'язаних зарядів діелектрика, оскільки вони лише коливаються відносно початкового
положення. Зміщення створює ефективний дипольний момент електрона $\vect p=-e\vect
r=-e^2\vect E/(m_e\omega^2)$, напрямлений проти поля, а~вектор поляризації одиниці об'єму
плазми
% ~~~~~~~~~~~~~~~~~~~~~~~~~~~~~~~~~~~~~~~~~~~~~~~~~~~~~~~~
\begin{equation*}
	\vect P=n_e\vect p=-\frac{n_ee^2}{m_e\omega^2}\vect E\equiv\chi\vect E.
\end{equation*}
% ~~~~~~~~~~~~~~~~~~~~~~~~~~~~~~~~~~~~~~~~~~~~~~~~~~~~~~~~
Оскільки $\vect D=\vect E+4\pi\vect P=\varepsilon\vect E$, діелектрична проникність
$\varepsilon=1+4\pi\chi$:
% ~~~~~~~~~~~~~~~~~~~~~~~~~~~~~~~~~~~~~~~~~~~~~~~~~~~~~~~~
\begin{equation}\label{eq:plasma-permittivity}
	\tcbhighmath{\varepsilon(\omega)=1-\frac{4\pi n_ee^2}{m_e\omega^2}=1-\frac{\omega_p^2}{\omega^2}}.
\end{equation}
% ~~~~~~~~~~~~~~~~~~~~~~~~~~~~~~~~~~~~~~~~~~~~~~~~~~~~~~~~~~
Той самий результат дає опис через струм. Швидкість електрона $\vect
v=\dot{\vect{r}}=-i\omega\vect r$, густина струму $\vect j=-n_ee\vect
v=i\dfrac{n_ee^2}{m_e\omega}\vect E\equiv\sigma(\omega)\vect E$, тому
% ~~~~~~~~~~~~~~~~~~~~~~~~~~~~~~~~~~~~~~~~~~~~~~~~~~~~~~~~
\begin{equation}\label{eq:sigma-collisionless}
	\sigma(\omega)=i\frac{\omega_p^2}{4\pi\omega},\qquad
	\varepsilon=1+\frac{4\pi i\sigma}{\omega}=1-\frac{\omega_p^2}{\omega^2}.
\end{equation}
% ~~~~~~~~~~~~~~~~~~~~~~~~~~~~~~~~~~~~~~~~~~~~~~~~~~~~~~~~~~~~
Провідність уявна: струм відстає від поля на чверть періоду (дія інерції електронів), тому
енергія не поглинається. Формула~\eqref{eq:plasma-permittivity}~--- окремий випадок моделі
Лоренца $\varepsilon=1+4\pi n_ee^2/[m_e(\omega_0^2-\omega^2)]$ при власній частоті
$\omega_0=0$: у~плазмі електрони \emph{вільні}, повертаючої сили немає.

Результат якісно відрізняється від електростатики діелектриків: тут $\varepsilon<1$, а~при
$\omega<\omega_p$ взагалі $\varepsilon<0$. Звернімо увагу ще на одне. Статичне екранування
(§\ref{sec:plasma-basics}) теж можна записати мовою діелектричної проникності: Фур'є-образ
потенціалу~\eqref{eq:screened-potential} дорівнює $\varphi_k=4\pi
q/(k^2+\lambda_D^{-2})=4\pi q/(k^2\varepsilon(k))$, де
$\varepsilon(0,k)=1+1/(k\lambda_D)^2$. Тобто статичне поле в~плазмі сильно перебудовується
через перерозподіл зарядів, а~високочастотне поле бачить уже іншу, інерційну відповідь
електронів. Це хороший приклад того, що \emphx{одне й~те саме середовище має різну
електричну відповідь на різних часових масштабах}.



%% --------------------------------------------------------
\section{Закон дисперсії. Частота відсікання}
%% --------------------------------------------------------



Для плоскої поперечної хвилі $\vect E=\vect E_0\exp(ikx-i\omega t)$ ($\vect k\cdot\vect
E=0$) рівняння Максвелла з~$\vect D=\varepsilon\vect E$ дають
$k^2c^2=\varepsilon\omega^2$. Підставляючи~\eqref{eq:plasma-permittivity}, одержуємо



\subsection*{Закон дисперсії}



% ~~~~~~~~~~~~~~~~~~~~~~~~~~~~~~~~~~~~~~~~~~~~~~~~~~~~~~~~
\begin{equation}\label{eq:plasma-dispersion}
	\tcbhighmath{\omega^2=\omega_p^2+c^2k^2}
\end{equation}
% ~~~~~~~~~~~~~~~~~~~~~~~~~~~~~~~~~~~~~~~~~~~~~~~~~~~~~~~~~~
(рис.~\ref{pic:plasma-dispersion}). Множачи на $\hbar^2$, бачимо, що це релятивістський
зв'язок $\mathcal E^2=(\hbar\omega_p)^2+(\hbar ck)^2$: у~плазмі фотон поводиться як
частинка з~\enquote{масою} $\hbar\omega_p/c^2$.

\begin{figure}[h!]\centering
	\localinput{dispersion.tikz}
	\caption{Закон дисперсії поперечної електромагнітної хвилі в~холодній плазмі. Частота
	$\omega_p$~--- границя поширення в~однорідній плазмі}
	\label{pic:plasma-dispersion}
\end{figure}

Хвилі з~$\omega>\omega_p$ поширюються з~довжиною $\lambda=2\pi
c/\sqrt{\omega^2-\omega_p^2}$, а~показник заломлення $n=\sqrt\varepsilon<1$. Швидкості
% ~~~~~~~~~~~~~~~~~~~~~~~~~~~~~~~~~~~~~~~~~~~~~~~~~~~~~~~~
\begin{equation*}
	v_{\mathrm{ф}}=\frac\omega k=\frac c{\sqrt{1-\omega_p^2/\omega^2}}>c,\qquad
	v_{\mathrm{гр}}=\frac{d\omega}{dk}=c\sqrt{1-\frac{\omega_p^2}{\omega^2}}<c,\qquad
	v_{\mathrm{ф}}v_{\mathrm{гр}}=c^2.
\end{equation*}
% ~~~~~~~~~~~~~~~~~~~~~~~~~~~~~~~~~~~~~~~~~~~~~~~~~~~~~~~~
Фазова швидкість більша за швидкість світла, і~це не суперечить теорії відносності, адже
переносять енергію й~сигнал хвильові пакети з~груповою швидкістю $v_{\mathrm{гр}}<c$.

\paragraph{Хвилі з~$\omega<\omega_p$.}
Тоді $k^2=(\omega^2-\omega_p^2)/c^2<0$. Покладемо $k=i\kappa$,
$\kappa=\sqrt{\omega_p^2-\omega^2}/c$:
% ~~~~~~~~~~~~~~~~~~~~~~~~~~~~~~~~~~~~~~~~~~~~~~~~~~~~~~~~
\begin{equation*}
	E=E_0\exp(-\kappa x-i\omega t).
\end{equation*}
% ~~~~~~~~~~~~~~~~~~~~~~~~~~~~~~~~~~~~~~~~~~~~~~~~~~~~~~~~
Поле експоненціально спадає в~глибину плазми, а~середній потік енергії дорівнює нулю.
Глибина проникнення $\delta=1/\kappa\approx c/\omega_p$ при $\omega\ll\omega_p$
(\emph{безколізійна скін-довжина}). Для $n_e=10^{12}~\mathrm{см}^{-3}$ вона становить
$\approx0{,}5$~см.

Отже, у~цій простій моделі однорідної, немагнітної та беззіткнювальної плазми поперечні
електромагнітні хвилі поширюються лише за $\omega>\omega_p$. Магнітне поле, градієнти
густини та зіткнення змінюють цю картину. Частоту $\omega_p$ називають \emphx{частотою
відсікання} (англ.~\emph{cutoff}).



%% --------------------------------------------------------
\section{Поздовжні (ленгмюрівські) хвилі}\label{sec:langmuir-waves}
%% --------------------------------------------------------



Запис $\vect D=\varepsilon\vect E$ і~рівняння $\nabla\cdot\vect D=0$ дають
$\varepsilon\,\vect k\cdot\vect E=0$. Є дві можливості: або $\vect k\cdot\vect E=0$
(поперечна хвиля, розглянута вище), або $\varepsilon(\omega)=0$ (поздовжня хвиля, $\vect
E\parallel\vect k$). З~\eqref{eq:plasma-permittivity} умова $\varepsilon=0$ дає
$\omega=\omega_p$~--- це відомі плазмові коливання (§\ref{sec:plasma-oscillations}).
У~холодній плазмі їхня частота не залежить від $k$, групова швидкість нульова, і~збурення
не поширюється.

Ситуація змінюється, якщо врахувати тепловий рух. Розглянемо одновимірні поздовжні хвилі
малої амплітуди ($n_e=n_0+n_1$, швидкість $v$, поле $E$; іони нерухомі):
% ~~~~~~~~~~~~~~~~~~~~~~~~~~~~~~~~~~~~~~~~~~~~~~~~~~~~~~~~
\begin{equation*}
	\frac{\partial n_1}{\partial t}+n_0\frac{\partial v}{\partial x}=0,\qquad
	m_en_0\frac{\partial v}{\partial t}=-en_0E-\frac{\partial p_1}{\partial x},\qquad
	\frac{\partial E}{\partial x}=-4\pi en_1.
\end{equation*}
% ~~~~~~~~~~~~~~~~~~~~~~~~~~~~~~~~~~~~~~~~~~~~~~~~~~~~~~~~
Стискання відбувається вздовж одного напрямку й~достатньо швидко, щоб було адіабатичним.
Для одного ступеня свободи показник адіабати $\gamma=(f+2)/f=3$, тобто $p\propto n^3$,
і~$p_1=3k_{\mathrm B}T_e\,n_1$. Виключаючи $v$ і~$E$, одержуємо
% ~~~~~~~~~~~~~~~~~~~~~~~~~~~~~~~~~~~~~~~~~~~~~~~~~~~~~~~~
\begin{equation}\label{eq:bohm-gross}
	\frac{\partial^2n_1}{\partial t^2}=-\omega_{pe}^2n_1+3v_{Te}^2\frac{\partial^2n_1}{\partial x^2}
	\quad\Rightarrow\quad
	\tcbhighmath{\omega^2=\omega_{pe}^2+3k^2v_{Te}^2}.
\end{equation}
% ~~~~~~~~~~~~~~~~~~~~~~~~~~~~~~~~~~~~~~~~~~~~~~~~~~~~~~~~
Це співвідношення Бома~--- Гросса. Воно є добрим прикладом узгодження кінетичного та
рідинного описів у~границі $k\lambda_{De}\ll1$. Тепер збурення поширюється з~груповою
швидкістю $v_{\mathrm{гр}}=3kv_{Te}^2/\omega\neq0$. Формула справедлива при
$k\lambda_{De}\ll1$. Коли $k\lambda_{De}$ перестає бути малим, рідинне наближення втрачає
точність, і~потрібний кінетичний опис. Зокрема, поблизу резонансу електрони, швидкість
яких близька до фазової швидкості хвилі, можуть забирати енергію хвилі навіть за
відсутності зіткнень (загасання Ландау, 1946~р.). Сила цього ефекту залежить від форми
розподілу та параметра $k\lambda_{De}$; це не просто поріг, після якого будь-яка хвиля
миттєво згасає. Це чисто кінетичний ефект, який неможливо одержати в~гідродинамічному
описі.




%% --------------------------------------------------------
\section{Зіткнення й~поглинання електромагнітних хвиль}\label{sec:wave-absorption}
%% --------------------------------------------------------

Досі зіткненнями ми нехтували. Врахуємо їх у~найпростішій формі (частоту $\nu\equiv\nu_{ei}$ обчислено в~розд.~\ref{sec:plasma-collisions}): електрон гальмується
силою тертя $-m_e\nu\vect v$, де $\nu$~--- ефективна частота зіткнень з~іонами та
нейтралами. Рівняння руху в~полі $\vect E\,e^{-i\omega t}$
% ~~~~~~~~~~~~~~~~~~~~~~~~~~~~~~~~~~~~~~~~~~~~~~~~~~~~~~~~
\begin{equation*}
	m_e\dot{\vect v}=-e\vect E-m_e\nu\vect v
	\quad\Rightarrow\quad
	\vect v=\frac{-ie\vect E}{m_e(\omega+i\nu)},
\end{equation*}
% ~~~~~~~~~~~~~~~~~~~~~~~~~~~~~~~~~~~~~~~~~~~~~~~~~~~~~~~~
звідки провідність $\sigma(\omega)=ie^2n_e/[m_e(\omega+i\nu)]$ і~діелектрична проникність
% ~~~~~~~~~~~~~~~~~~~~~~~~~~~~~~~~~~~~~~~~~~~~~~~~~~~~~~~~
\begin{equation}\label{eq:eps-collisions-waves}
	\varepsilon(\omega)=1+\frac{4\pi i\sigma}{\omega}=1-\frac{\omega_p^2}{\omega(\omega+i\nu)}.
\end{equation}
% ~~~~~~~~~~~~~~~~~~~~~~~~~~~~~~~~~~~~~~~~~~~~~~~~~~~~~~~~
Тепер $\varepsilon$ комплексна, отже, і~хвильове число $k=(\omega/c)\sqrt\varepsilon$ має
уявну частину: хвиля поглинається. Є дві граничні ситуації.
\begin{itemize}
	\item $\omega\gg\nu$ (високочастотна границя). Тоді
	  $\operatorname{Im}\varepsilon\approx\omega_p^2\nu/\omega^3$, і~амплітуда спадає як
	  $e^{-\operatorname{Im}k\,x}$ з
	  \begin{equation*}
	  	\operatorname{Im}k\approx\frac{\nu}{2c}\,\frac{\omega_p^2/\omega^2}{\sqrt{1-\omega_p^2/\omega^2}}.
	  \end{equation*}
	  Поглинання пропорційне $\nu$ і~різко зростає поблизу відсікання $\omega\to\omega_p$.
	  Фізично: коливний рух електрона тим сильніше переходить у~тепло, чим частіше він
	  зіштовхується. Цей механізм поглинає радіохвилі в~шарі~D іоносфери.
	\item $\omega\ll\nu$ (низькочастотна границя). Електрони рухаються як у~постійному
	  полі: $\sigma\to\sigma_0=\omega_p^2/4\pi\nu$, плазма поводиться як звичайний провідник
	  з~товщиною скін-шару $\delta=c/\sqrt{2\pi\sigma_0\omega}$. Вона зростає з~нагріванням
	  плазми, бо $\nu\propto T^{-3/2}$.
\end{itemize}
Отже, \emphx{чи є плазма поглинальною, вирішує порівняння $\omega$ з~$\nu$}, а~чи
відбивальною~--- порівняння $\omega$ з~$\omega_p$. Джоулеве тепло $\langle\vect j\cdot\vect
E\rangle=\operatorname{Re}\sigma\,|E|^2/2$ використовують для нагрівання плазми
НВЧ-полем.

%% --------------------------------------------------------
\section{Відбиття електромагнітних хвиль від плазми}
%% --------------------------------------------------------



Нехай на межу плазми ($x>0$) нормально падає хвиля з~вакууму. Амплітудний коефіцієнт
відбиття (формули Френеля, $\mu=1$)
% ~~~~~~~~~~~~~~~~~~~~~~~~~~~~~~~~~~~~~~~~~~~~~~~~~~~~~~~~
\begin{equation*}
	r=\frac{1-\sqrt\varepsilon}{1+\sqrt\varepsilon}.
\end{equation*}
% ~~~~~~~~~~~~~~~~~~~~~~~~~~~~~~~~~~~~~~~~~~~~~~~~~~~~~~~~
При $\omega<\omega_p$ величина $\sqrt\varepsilon=i\kappa$ чисто уявна, тому
% ~~~~~~~~~~~~~~~~~~~~~~~~~~~~~~~~~~~~~~~~~~~~~~~~~~~~~~~~
\begin{equation*}
	R=|r|^2=\left|\frac{1-i\kappa}{1+i\kappa}\right|^2=1.
\end{equation*}
% ~~~~~~~~~~~~~~~~~~~~~~~~~~~~~~~~~~~~~~~~~~~~~~~~~~~~~~~~
При $\omega\gg\omega_p$ показник заломлення наближається до одиниці, $n\to1$, а~коефіцієнт
відбиття прямує до нуля, $R\to0$. Плазма стає прозорою для електромагнітного
випромінювання.

Саме тому високочастотне випромінювання, для якого $\omega\gg\omega_p$, може проходити
крізь плазму; до такого діапазону, зокрема, належить рентгенівське випромінювання для
багатьох плазмових середовищ.



\subsection*{Похиле падіння}



При $\omega>\omega_p$ плазма~--- середовище з~дійсним показником заломлення
$n=\sqrt\varepsilon<1$, тобто вакуум для хвилі оптично густіший. Тому можливе повне
внутрішнє відбиття при $\sin\theta\geq n$. Критичний кут
$\sin\theta_{\mathrm{кр}}=\sqrt{1-\omega_p^2/\omega^2}$ прямує до нуля при
$\omega\to\omega_p+0$, і~плазма відбиває при будь-якому куті падіння.



\subsection*{Плавна межа}



Реальні плазми (іоносфера) мають концентрацію, що повільно зростає з~глибиною:
$n_e=n_e(x)$ на відстанях порядку довжини хвилі змінюється слабко. Тому можна використати
локальне значення хвильового числа
% ~~~~~~~~~~~~~~~~~~~~~~~~~~~~~~~~~~~~~~~~~~~~~~~~~~~~~~~~
\begin{equation*}
	k^2(x)=\frac1{c^2}\left[\omega^2-\frac{4\pi e^2}{m_e}n_e(x)\right].
\end{equation*}
% ~~~~~~~~~~~~~~~~~~~~~~~~~~~~~~~~~~~~~~~~~~~~~~~~~~~~~~~~
При нормальному падінні $k(x)$ перетворюється на нуль, де концентрація досягає
\emphx{критичного значення}
% ~~~~~~~~~~~~~~~~~~~~~~~~~~~~~~~~~~~~~~~~~~~~~~~~~~~~~~~~
\begin{equation}\label{eq:critical-density}
	n_{\mathrm{кр}}=\frac{m_e\omega^2}{4\pi e^2}\approx1{,}24\times10^{4}\,f^2[\mathrm{МГц}]~\mathrm{см}^{-3},\qquad f=\frac\omega{2\pi}.
\end{equation}
% ~~~~~~~~~~~~~~~~~~~~~~~~~~~~~~~~~~~~~~~~~~~~~~~~~~~~~~~~~~~~~
У~цій точці (\enquote{плазмове дзеркало}) хвиля повертається назад. Строго кажучи,
наближення локального $k(x)$ у~самій точці повороту непридатне (потрібне розв'язання
рівняння Ейрі), але положення точки воно визначає правильно.

При похилому падінні під кутом $\theta_0$ з~вакууму в~шарувате середовище зберігається
тангенційна складова хвильового вектора, тобто закон Снелліуса
$\sqrt{\varepsilon(x)}\sin\theta(x)=\sin\theta_0$. Поворот відбувається там, де
$\theta=90^\circ$:
% ~~~~~~~~~~~~~~~~~~~~~~~~~~~~~~~~~~~~~~~~~~~~~~~~~~~~~~~~
\begin{equation*}
	\varepsilon(x_t)=\sin^2\theta_0\quad\Rightarrow\quad n_e(x_t)=n_{\mathrm{кр}}\cos^2\theta_0.
\end{equation*}
% ~~~~~~~~~~~~~~~~~~~~~~~~~~~~~~~~~~~~~~~~~~~~~~~~~~~~~~~~
Тому хвилю, що падає похило, плазма відбиває за меншої концентрації: максимальна частота,
яку може відбити шар з~максимальною концентрацією $n_{\max}$, дорівнює
$f_{\max}=f_p(n_{\max})/\cos\theta_0$ (\enquote{закон секанса}).

% \begin{figure}[h!]\centering
% \localinput{gradient-plasma.tikz}
% 	\caption{Хвиля падає на плазму зі зростаючою концентрацією електронів: точка повороту
% 	розташована там, де $n_e=n_{\text{кр}}\cos^2\theta_0$}
% 	\label{pic:gradient-plasma}
% \end{figure}



%% --------------------------------------------------------
\section{Іоносфера і~далекий радіозв'язок}
%% --------------------------------------------------------



Іоносфера~--- шар верхньої атмосфери (від висоти $\sim60$~км), у~якому газ іонізований
ультрафіолетовим і~рентгенівським випромінюванням Сонця та космічними променями. Ступінь
іонізації невелика, але концентрація електронів (від $\sim10^{4}$~$\mathrm{см}^{-3}$
в~нижніх областях до $\sim10^{6}~\mathrm{см}^{-3}$ у~максимумі шару~F на висотах
$200$--$400$~км вдень) достатня для відбиття радіохвиль. Для $n_e=10^{6}~\mathrm{см}^{-3}$
частота $f_p\approx9$~МГц, тобто короткі хвилі (декаметрові) відбиваються, а~ультракороткі
й~мікрохвилі проходять крізь іоносферу~--- тому ними користуються для зв'язку з~космічними
апаратами.

Відбиття від іоносфери дозволяє приймати радіосигнали за горизонтом
(рис.~\ref{pic:radio-communication}). Довгі та середні хвилі поводяться інакше: вдень їх
поглинає нижній шар~D через часті зіткнення електронів з~нейтральними молекулами
(§\ref{sec:plasma-collisions}), а~вночі шар~D зникає, і~сигнали відбиваються від вищих
шарів. Звідси відомий факт: вночі середньохвильові радіостанції чути значно далі.

\begin{figure}[h!]\centering
	% \figstub[0.75]{Земля, іоносфера, ламана траєкторія променя між передавачем
	% і~приймачем}
	\includegraphics[width=0.75\linewidth]{radio-communication}
	\caption{Далекий радіозв'язок завдяки відбиттю від іоносфери}
	\label{pic:radio-communication}
\end{figure}

Той самий ефект проявляється в~інших ситуаціях. Під час входу космічного апарата
в~атмосферу навколо нього утворюється плазмова оболонка з~концентрацією
$n_e\gtrsim10^{11}~\mathrm{см}^{-3}$, що за формулою~\eqref{eq:critical-density} відсікає
радіозв'язок на частотах нижче кількох гігагерц (\enquote{радіозатемнення}). Метали мають
вільні електрони з~$n_e\sim10^{23}~\mathrm{см}^{-3}$ і~$\hbar\omega_p\sim10$~еВ
(ультрафіолет): вони відбивають видиме світло за тим самим механізмом, що й~плазма
радіохвилі, а~в~ультрафіолеті для багатьох металів (наприклад, лужних) стають прозорими.





%% --------------------------------------------------------
%% --------------------------------------------------------
\begin{Summary}
\begin{itemize}
\item Плазмова частота $\omega_p$~--- частота власних коливань електронів відносно іонів; вона задає часовий масштаб колективного відгуку.
\item Діелектрична проникність холодної плазми $\varepsilon=1-\omega_p^2/\omega^2$; поперечні хвилі поширюються лише за $\omega>\omega_p$.
\item При $\omega<\omega_p$ хвиля відбивається: так працюють іоносфера й~далекий радіозв'язок.
\item Поздовжні ленгмюрівські хвилі мають тепловий поправковий член у~законі дисперсії; зіткнення призводять до поглинання хвиль.
\end{itemize}
\end{Summary}

\Needspace{9\baselineskip}
\section*{Задачі}
%% --------------------------------------------------------

\begin{problem}\label{prb:PlasmaFrequency}
Знайдіть плазмову частоту $\omega_{pe}$, частоту $f=\omega_{pe}/2\pi$ та період коливань
для $n_e=10^{10}~\mathrm{см}^{-3}$. До якого діапазону електромагнітних хвиль належить ця
частота?
\end{problem}

\begin{problem}\label{prb:PlasmaIons}
Покажіть, що відношення іонного та електронного внесків у~$\omega_p^2$ для водневої плазми
дорівнює $m_e/m_i$, і~оцініть відносну поправку до плазмової частоти.
\end{problem}

\begin{problem}\label{prb:BohmGross}
Для електронної плазми з~температурою $T_e$ оцініть поправку $3k^2v_{Te}^2$ до квадрата
ленгмюрівської частоти. За яких умов ця поправка мала порівняно з~$\omega_{pe}^2$?
\end{problem}

\begin{problem}\label{prb:PlasmaWave}
	Радіохвиля частоти 1~МГц падає на область плазми з~концентрацією
	$n_e=10^6~\mathrm{см}^{-3}$. Чи проникне хвиля в~плазму? Оцініть глибину проникнення.
\end{problem}

\begin{problem}\label{prb:PlasmaIonosphere}
	Максимальна концентрація шару~F дорівнює $10^6~\mathrm{см}^{-3}$. Яку максимальну
	частоту хвилі може відбити іоносфера при кутах падіння $0^\circ$ і~$70^\circ$? Чому
	сигнал 100~МГц практично не відбивається?
\end{problem}

%% --------------------------------------------------------
%% --------------------------------------------------------

\begin{Questions}
\begin{quizlist}
\item Чому плазма відбиває хвилі з~$\omega<\omega_p$ повністю? Що відбувається при
		  $\omega>\omega_p$?

\item Як пояснити далекий радіозв'язок на коротких хвилях? Чому ультракороткі
		  хвилі проходять крізь іоносферу?

\item Чому при поширенні в~плазмі з'являється поглинання? Що вирішує:
		  $\omega\gg\nu$ чи $\omega\ll\nu$? Чому поглинання зростає поблизу $\omega_p$?

\item Поясніть фізичну природу плазмових коливань. Чому частота не залежить від
		  $k$ у~холодній плазмі?

\item Як з~діелектричної проникності плазми одержати закон дисперсії? Чому
		  $v_{\mathrm{ф}}>c$ не порушує теорії відносності?

\end{quizlist}
\end{Questions}


